\section{Induced maps on homology}
\label{sec:defmaps}
In this section we define $\HH_\blob(X)$, the Hochschild homology of a
space $X$, and show that given a kernel $\Phi\colon X\kerto Y$ we get
pull-back and push-forward maps, $\Phi^*\colon \HH_\blob(Y)\to
\HH_\blob(X)$ and $\Phi_*\colon \HH_\blob(X)\to \HH_\blob(Y)$, such
that if $\Phi$ is right adjoint to $\Psi$ then $\Psi^*=\Phi_*$.

 \subsection{Hochschild cohomology.}
First recall that for a space $X$, one way to define its Hochschild
\emph{cohomology} is as
\[\HH^\blob(X):=\Ext_{X\times X}^\blob(\OO_\Delta,\OO_\Delta).\]
However, the ext-group is just the hom-set $\Hom^\blob_{\D(X\times
  X)}(\OO_\Delta,\OO_\Delta)$ which by the definition of $\gVar$ is
just $\THom^\blob_\gVar(\Id_X, \Id_X)$.  In terms of diagrams, we can
thus denote an element $\phi\in \HH^\blob(X)$ as
\[\pic{phiinHH}.\]
Note that the grading is not indicated in the picture, but this should
not give rise to confusion.

 \subsection{Hochschild homology.}
Now we define $\HH_\blob(X)$ the Hochschild \emph{homology} of a space
$X$ as follows:
\[\HH_\blob(X):=\THom_\gVar^\blob(\SerK_X^{-1},\Id_X) \]
or, in more concrete terms, 
\[ \HH_\blob(X) = \Ext_{X\times X}^{-\blob}(\SerK^{-1}_X,\OO_\Delta).\]
Thus an element $w\in \HH_\blob(X)$ will be denoted
\[\pic{winHH},\]
where again the shifts are understood.


It is worth taking a moment to compare this with other definitions of
Hochschild homology, such as that of Weibel~\cite{Weibel}.  He defines
the Hochschild homology of a space $X$ as
$\H^\blob(X,\Delta^*\cO_\Delta)$, where as usual by $\Delta^*$ we mean
the left-derived functor.  This cohomology group is naturally
identified with the hom-set $\Hom^\blob_{X}(\cO_X,\Delta^*\cO_\Delta)$
which is isomorphic to $\Hom^\blob_{X\times
  X}(\Delta_!\cO_X,\cO_\Delta)$ where $\Delta_!$ is the left-adjoint
of $\Delta^*$.  Direct calculation shows that $\Delta_!\cO_X\cong
\SerK_X^{-1}$ and so our definition is recovered.  Another feasible
definition of Hochschild homology is $\H^\blob(X\times
X,\cO_\Delta\otimes \cO_\Delta)$, and this again is equivalent to our
definition as there is the isomorphism $\SerK_X^{-1}\cong
\cO_\Delta^\chk$.

 \subsection{Push-forward and pull-back.}
\label{def:defmaps}
For spaces $X$ and $Y$ and a kernel $\Phi\colon X\kerto Y$ define the
push-forward on Hochschild homology $\Phi_*\colon \HH_\blob(X)\to \HH_\blob(Y)$ as follows
\[\Phi_*\left(\pic{winHH}\right):=\pic{PhiPushOfw}.\]
For the reader still unhappy with diagrams, for $v\in
\Hom^\blob(\SerK_X^{-1},\Id_X)$, define
$\Phi_*(v)\in \Hom^\blob(\SerK_Y^{-1},\Id_Y)$ as
the following composite, which is read from the above diagram by
reading upwards from the bottom:
\begin{align*}
  \SerK^{-1}_Y
  &\xrightarrow{\gamma}
  \Phi\circ \Phi^\chk
  \xrightarrow{\Id\circ\eta\circ\Id}
  \Phi\circ\SerK_X^{-1}\circ\SerK_X\circ \Phi^\chk
  \xrightarrow{\Id\circ v\circ\Id\circ\Id} \Phi\circ\SerK_X\circ
  \Phi^\chk \xrightarrow{\rote} \Id_Y.
\end{align*}
\medskip

\noindent
Similarly define the pull-back $\Phi^*\colon \HH_\blob(Y)\to
\HH_\blob(X)$ as follows
\[\Phi^*\left(\pic{vinHH}\right):=\pic{PhiPullOfv}.\]
These operations depend only on the isomorphism class of the kernel as
shown by the following.

\begin{Proposition}
\label{prop:indepiso}
  If kernels $\Phi$ and $\hat\Phi$ are isomorphic then they give rise to
  equal push-forwards and equal pull-backs: $\Phi_*=\hat\Phi_*$ and
  $\Phi^*=\hat\Phi^*$.
\end{Proposition}

\begin{Proof}
  This follows immediately from the fact that the 2-morphisms $\gamma_\Phi$,
  $\gamma_{\hat\Phi}$, $\rote_\Phi$ and $\rote_{\hat\Phi}$ of
  Section~\ref{subsec:yphimphi} are natural and thus commute with the
  given kernel isomorphism $\Phi\cong\hat\Phi$.
\end{Proof}
\medskip

\noindent
The push-forward and pull-back operations are functorial in the
following sense.

\begin{Theorem}[Functoriality]
\label{thm:functorial}
  If $\Phi\colon X\to Y$ and $\Psi\colon Y\to Z$ are kernels then the
  push-forwards and pull-backs compose appropriately, namely:
  \[(\Psi\circ\Phi)_*=\Psi_*\circ \Phi_*\colon
  \HH_\blob(X)\to\HH_\blob(Z)\]
  and
  \[(\Psi\circ\Phi)^*=\Phi^*\circ \Psi^*\colon
  \HH_\blob(Z)\to\HH_\blob(X).\]
\end{Theorem}

\begin{Proof}
  This follows from the fact that the right adjunction $\tau_R$ is a
  2-functor.  The adjoint of $\Psi\circ\Phi$ is canonically
  $\tau_R(\Phi)\circ\tau_R(\Psi)$, i.e., is
  $\SerK_X\circ\Phi^\chk\circ\SerK_Y\circ\Psi^\chk$.  This means that
  the unit of the adjunction $\Id\Rightarrow
  \Psi\circ\Phi\circ\SerK_X\circ\Phi^\chk\circ\SerK_Y\circ\Psi^\chk$
  is given by the composition
  $\Id\Rightarrow\Psi\circ\SerK_Y\circ\Psi^\chk\Rightarrow
  \Psi\circ\Phi\circ\SerK_X\circ\Phi^\chk\circ\SerK_Y\circ\Psi^\chk$.
  This gives
  \[(\Psi\circ\Phi)_*(w)=\pic{PsiOfPhiOfw}=\Psi_*\left(\Phi_*(w)\right).
  \qedhere
  \]
\end{Proof}

\begin{Theorem}
\label{thm:adjadj}
  If $\Phi\colon X\to Y$ and $\Psi\colon Y\to X$ are adjoint kernels,
  $\Phi\dashv\Psi$, then we have
  \[\Phi_*=\Psi^*\colon \HH_\blob(X)\to\HH_\blob(Y).\]
\end{Theorem}

\begin{Proof}
  By the uniqueness of adjoints we have a canonical isomorphism,
  $\Psi\cong\tau_R(\Phi)$ and by Proposition~\ref{prop:indepiso}
  $\Psi^*=(\tau_R(\Phi))^*$.  It therefore suffices to show that
  $\Phi_*=(\tau_R(\Phi))^*$.  Observe that
  \[\tau_R(\tau_R(\Phi))=\tau_R(\SerK_X\circ\Phi^\chk)\cong
  \tau_R(\Phi^\chk)\circ
  \tau_R(\SerK_X)\cong\SerK_Y\circ\Phi^{\chk\chk}\circ \SerK_X^{-1},\]
  and similarly
 \[\tau_L(\tau_R(\Phi))\cong
 \Phi\circ\SerK_X^{-1}\circ \SerK_X.\] Of course the latter is
 isomorphic to $\Phi$ but the Serre kernels are left in so to make the
 adjunctions more transparent.  We now get the unit for adjunction
 $\tau_R(\Phi)\dashv\tau_R(\tau_R(\Phi))$ and the counit for the
 adjunction $\tau_L(\tau_R(\Phi))\dashv\tau_R(\Phi)$ as follows:
 \[\pic{tauRunit};\quad\pic{tauRcounit}.\]
 Thus
 \[\tau_R(\Phi)^*\left(w\right)
 =\pic{MessyPhiPushOfw}= \pic{PhiPushOfw}=\Phi_*(w). \qedhere\]
\end{Proof}
